\heading{数据结构与算法A-15秋-期末}
\noteinfo{题目来源为真题扫描件转录。客观题答案和部分参考解答已整理在解答中；原扫描中考场纪律、装订线等非题目信息已尽量删去。本次整理提供 C 语言版与 Python 语言版；除程序代码语言外，两版题目内容保持一致。}

\begin{question}
一、选择题（每题 2 分，共 20 分）\\
1.\\
若要一个运行时间为$100n^2$ 的算法在相同机器上快于另一个运行时间为2n 的算法，则最小\\
的n 值应为\_\_\_\_\_\_\_\_\_。\\
2.\\
某算法仅含顺序执行的语句1 和语句2，语句1 执行次数为$20n^2$+3nlogn，而语句2 执行次\\
数为$2n^3$，则该算法的时间复杂度为 \_\_\_\_\_\_\_\_\_\_\_\_\_\_\_。\\
3.\\
下面术语中， \_\_\_\_\_\_\_\_\_\_\_\_\_与数据结构的存储结构无关。\\
A． 顺序表       B.  链表      C.  队列     D. 循环链表     E. 堆\\
4.\\
在一个具有n 个结点的有序单链表中插入一个新结点并仍保持其有序的时间复杂度为\\
\_\_\_\_\_\_\_\_\_\_\_\_。\\
A.  $O(n\log_2 n)$      B.  O(1)       C.  O(n)\\
D.  $O(n^2)$\\
5.\\
若元素a、b、c、d、e、f 依次进栈，允许进栈、退栈操作交替进行，但不允许连续三次进\\
行退栈操作，则不可能得到的出栈序列是 \_\_\_\_\_\_\_\_\_。 （\quad）\\
A.  dcebfa\\
B.  cbdaef\\
C.  bcaefd\\
D.  afedcb\\
6.\\
表达式3* 2\textasciicircum{}(4+2*2-6*3)-5 的求值过程中，当扫描到6 时，操作数栈和运算符栈为       ，\\
其中 \textasciicircum{} 为乘幂。 （\quad）\\
A.  3,2,4,1,1；(*\textasciicircum{}(+*-\\
B.  3,2,8；(*\textasciicircum{}-\\
C.  3,2,4,2,2；(*\textasciicircum{}(-\\
D.  3,2,8；(*\textasciicircum{}(-\\
7.\\
有4 棵结点存储整数关键码的二叉树，若它们的中序遍历序列分别为：   A.  5, 3, 1, 2, 4\\
B.  1, 2, 3, 4, 5    C.   5, 4, 3, 2, 1   D.   1, 3, 5, 4, 2。请问其中可能是二叉搜索树的包括\\
\_\_\_\_\_\_\_\_\_\_。\\
8.\\
（4 分）一棵Huffman 树的带权外部路径长度定义为所有叶结点权值与其深度的乘积之和。\\
设根结点的深度为0，对集合 \{ 2, 3, 5, 7, 8 \} 构造二叉Huffman 树，则其最小带权外部路径\\
长度为\_\_\_\_\_\_\_\_\_\_\_\_\_\_\_。\\
9.\\
两个字符串相等的充分必要条件是   \_\_\_\_\_\_\_\_\_\_\_\_\_\_\_\_\_\_\_\_\_\_\_\_\_\_\_\_\_\_\_\_\_\_\_\_\_\_\_。(两字\\
符串的长度相等且两字符串中对应位置的字符也相等)\\
10. 2-3 树是一种满足以下两个条件的特殊的树： （1）每个内部结点有两个或三个子结点；（2）\\
所有叶结点到根的路径长度相同。如果一棵2-3 树有9 个叶结点，那么它可能有\_\_\_\_\_\_\_\_\_\\
个非叶结点。\\
二、辨析与简答题（共 36 分）\\
1.\\
（6 分）什么是循环队列？循环队列的优点是什么？如何判别它的空和满？\\
2.\\
（6 分）在包含 n 个结点的单链表、双链表和单循环链表中，若仅知道指针p 指向某结点，\\
不知道表的头指针，能否将结点p 从相应的链表中删去? 若可以，其时间复杂度各为多少?\\
3.\\
（8 分）请分别计算模式p = “aabaac” 优化前和优化后的next 数组，并按照优化后的next\\
数组针对目标 t = "aabaabaabaac" 进行 KMP 快速模式匹配，请画出匹配过程的示意图并\\
计算匹配过程中的比较次数。\\
4.\\
（8 分）根据树的双标记层次遍历序列，求其带度数的后根遍历序列 。譬如，已知一棵树\\
的双标记层次遍历序列如下：\\
A : ltag: 0 , rtag 1\\
B : ltag: 0 , rtag 0\\
C : ltag: 0 , rtag 1\\
D : ltag: 1 , rtag 0\\
G : ltag: 0 , rtag 1\\
E : ltag: 0 , rtag 1\\
H : ltag: 1 , rtag 1\\
F : ltag: 1 , rtag 0\\
I : ltag: 1 , rtag 1\\
则其带度数的后根遍历序列为：D0 H0 G1 B2 F0 I0 E2 C1 A2 （注：各个结点按照 “结点名 度\\
数” 的方式给出，结点之间用空格分隔）。 现某棵树的双标记层次遍历序列如下：\\
A : ltag: 0 , rtag 1\\
B : ltag: 0 , rtag 0\\
G : ltag: 1 , rtag 1\\
C : ltag: 0 , rtag 0\\
F : ltag: 0 , rtag 1\\
D : ltag: 1 , rtag 0\\
E : ltag: 1 , rtag 1\\
H : ltag: 1 , rtag 0\\
I : ltag: 1 , rtag 1\\
请给出其带度数的后根遍历序列。\\
5.\\
（8 分）使用重量权衡合并规则（权重合并规则）与路径压缩，对下列从0 到15 之间的数\\
的等价对进行归并。在初始情况下，集合中的每个元素分别在独立的等价类中。当两棵树\\
规模同样大时，使结点数值较大的根结点作为值较小的根结点的子结点。\\
(0,2) (1,2) (3,4) (3,1) (3,5) (9,11) (12,14) (3,9) (4,14) (6,7) (8,10) (8,7) (7,0) (10,15) (10,13)\\
请填写下面表格的空白部分树的父指针表示法的数组表示。也就是所有等价对都被处理之\\
后，所得父结点的下标值（没有父结点则填“–1”）。\\
父结点下标\\
结点值\\
0 1\\
2 3\\
4 5\\
7 8\\
9 10 11 12 13 14 15\\
结点的下标 0 1\\
2 3\\
4 5\\
7 8\\
9 10 11 12 13 14 15\\
三、算法填空题（每题 12 分，共 24 分）\\
1.\\
（每空3 分，共12 分）下面的算法利用一个栈将一给定的序列从小到大排序。长度为n\\
的序列由1, 2, …, n 这n 个连续的正整数组成。请补全下面的代码段，使其可以判断给定的\\
序列是否可以利用一个栈进行排序，如果可以，输出相应的操作。\\
例如：序列4 3 1 2，此序列可以排序，输出：push push push pop push pop pop pop。\\
template <class T>                            // 栈的元素类型为T\\
class Stack \{\\
public:\\
void clear();\\
// 变为空栈\\
void push(const T item);\\
//  item 入栈\\
T pop();\\
// 返回栈顶内容并弹出\\
T top();\\
// 返回栈顶内容但不弹出\\
bool isEmpty();\\
// 若栈已空返回真\\
bool isFull();\\
// 若栈已满返回真\\
\};\\
Stack<int> s;\\
//  s 为栈\\
int sequence[maxLen],  len;                 // sequence 为待排序序列， len 为序列长度\\
bool islegal()  \{\\
// 判断序列是否可以排序\\
int i, j, k;\\
for (i = 0; i < len; i++) \{\\
for (j = i+1; j < len; j++) \{\\
for (k = j+1; k < len; k++) \{\\
if (            )\\
return false;\\
\}\\
\}\\
\}\\
return true;\\
\}\\
void transform()  \{\\
// 输出转换操作\\
int cur = 1,  i = 0;\\
while  (           ) \{\\
while (s.isEmpty()  ||            ) \{\\
cout << "push ";\\
\}\\
s.pop();\\
cout << "pop ";\\
++i;\\
\}\\
\}\\
2.\\
（每空3 分，分析 3 分，共12 分）以下算法用来判断两棵树是否相同，试补全下面的代\\
码段，并分析算法的运行时间代价。\\
template <class T>\\
bool IsEqual(TreeNode<T>  *root1,  TreeNode<T>  *root2) \{\\
while (root1 != NULL  \&\& root2 != NULL) \{\\
if (root1->value()  !=  root2->value())\\
return   \_\_\_\_\_\_ ;\\
if  (\_\_\_\_\_\_\_\_\_\_\_\_\_\_\_\_\_)\\
return false;\\
root1 = root1->RightSibling();\\
root2 = root2->RightSibling();\\
\}\\
if (\_\_\_\_\_\_\_\_\_\_\_\_\_\_)\\
return true;\\
return false;\\
\}\\
四、\\
分析证明题 (共 20 分)\\
1.\\
（5 分） 请证明非空满K 叉树的叶结点数目为 (K-1)*n+1，其中n 为分支结点数目。\\
2.\\
（15 分）有n 个数K1, K2, …, Kn 顺序存储在数组中，形成一棵n 个结点的完全二叉树。现在\\
要将它按如下方式建成一个最小值堆：从左至右扫描整个数组，将第i 个元素Ki (i = 1, 2, …,\\
n) 与其父结点K[i/2] 比较，如果Ki >= K[i/2]，则不再进行操作；若Ki < K[i/2]，将2 个结点对换，\\
再将K[i/2]与K[i/4] 比较，以此类推，直到比较到根结点K1。\\
(1) 按照这种建堆方式，最坏情况下建堆的时间复杂度是多少？\\
(2) 教材上使用的“筛选法”建堆 (Siftdown)，最坏情况下建堆的时间复杂度是多少？\\
(3) 上述2 种建堆法有类似之处，请问它们的时间复杂度是否一样，若不同，则请说明导\\
致它们时间复杂度不同的原因。\\
\begin{solution}
本题为整卷转录型资料：选择题、填空题的参考答案以及简答、算法设计题的参考做法已统一放在本解答中，题面不保留原扫描夹带的答案标注。对扫描不清处，保留了原转录中可辨认的信息，后续可继续按原 PDF 图示补充图片。

\medskip
\noindent 原转录中夹带在题面里的参考答案/解答摘录如下，现已移入解答：
本卷包含选择填空、简答、算法设计等题型。下面保留按扫描件整理的题面；其中客观题参考答案已填入空格或括号，主观题给出参考思路或代码。

参考答案：（建议每一问2分） （1）用一维数组表示队列，由于队列的性质（队尾插入和队头\\
删除），容易造成“假溢出”现象，即队尾已到达一维数组的最高下标，不能再插入，然而队中\\
元素个数小于队列的长度（容量）。循环队列是解决“假溢出”的一种方法。通常把一维数组看\\
成首尾相接。\\
（2）优点：可以克服顺序队列的"假上溢"现象，能够使存储队列的向量空间得到充分的利用；\\
（3）判别"空"或"满"方法：少用一个元素的空间，如果头尾指针相等则队列为空，如果入队前\\
测试入队后头尾指针重合，则队列已满。\\
参考答案：分链表的三种情况讨论：\\
a)\\
（２分）单链表。若指针p 指向某结点时，能够根据该指针p 找到其直接后继，且按\\
照后继指针链找到p 结点后的所有结点。但是由于不知道其头指针，所以无法访问p\\
指针指向结点的直接前驱。因此无法删去该结点。另一种思路（有若干同学采用）：可\\
将p 的后继结点的值赋给p, 并将其删除，O(1)的时间复杂度，虽然不是删除p 结点，\\
但从某种意义上可达到同样的效果；可酌情给分？   (不可删除给２分，采用后一种\\
者酌情可给１分)\\
b)\\
（２分）双链表。由于双链表提供双向指针，根据p 结点的前驱指针和后继指针可以\\
查找到其直接前驱和直接后继，从而可以删除该结点。其时间复杂度为O（1） 。（2\\
分，可否1 分，时间复杂度1 分）\\
c)\\
（２分）单循环链表。根据已知结点位置，可以直接得到其后相邻的结点位置，因是\\
循环链表，所以可以循链得到p 结点的直接前驱。因此可以删去p 所指结点。其时间\\
复杂度应为O(n). （2 分，可否1 分，时间复杂度1 分）\\
参考答案\\
p = “aabaac”优化后的特征向量为：\\
下标\\
P\\
a\\
a\\
b\\
a\\
a\\
c\\
N\\
-1\\
-1\\
-1\\
-1\\
a\\
a\\
b\\
a\\
a\\
b\\
a\\
a\\
b\\
a\\
a\\
c\\
a\\
a\\
b\\
a\\
a\\
c\\
×\\
i = 5 j = 5 i = n[i] = 2\\
a\\
a\\
b\\
a\\
a\\
c\\
×\\
i = 5 j = 8 i = n[i] = 2\\
a\\
a\\
b\\
a\\
a\\
c\\
匹配\\
比较次数为14\\
参考答案 ：D0 E0 C2 H0 I0 F2 B2 G0 A2\\
参考答案 ：\\
父结点下标 -1 0  0\\
0 0\\
0 0\\
0 0\\
12 0\\
结点值\\
3 4\\
5 6\\
8 9\\
10 11 12 13 14 15\\
结点的下标 0\\
3 4\\
5 6\\
8 9\\
10 11 12 13 14 15\\
参考答案：\\
(1) sequence[k] < sequence[i] \&\& sequence[i] < sequence[j]\\
(2) cur <= len\\
(3) cur < s.top()\\
(4) s.push(sequence[i++]);\\
参考答案：\\
(1)false\\
(2) false==IsEqual(root1->LeftMostChild(),root2->LeftMostChild()\\
(3) root1==NULL \&\& root2==NULL\\
由于每个结点只需要访问一次,则复杂度为O(n)。\\
参考答案：\\
正确。数学归纳法证明如下：\\
1) 当n=0 时，只有一个根结点，此时叶结点数目为1，而(K-1)n+1 也为1\\
n=1 时，(K-1)n+1=K, 这是只有一个分支结点的情况，结果正确；\\
2) ．假设n=m ( m>=1, m 为自然数)时叶结点数目为 （K-1）m+1,\\
则n=m+1 时，相当于把n=m 时的K 叉树里的一个叶结点扩展为分支结点，叶结点增加\\
K-1 个\\
故n=m+1 时，叶结点数目为(K-1)m+1+K-1=(m+1)K-m\\
=(m+1)K-(m+1)+1\\
=(K-1)(m+1)+1\\
=(K-1)n+1\\
即n=m+1 时此公式仍然成立。\\
综1、2 所述，该公式成立，证毕。\\
参考答案：\\
(1)最坏情况下，树上所有的元素都要比较至根结点。所以第i层(0<=i<=logn，\\
根为第0层)的结点均比较了i次。而第i层有2i个结点，因此总的比较次数为\\
从而S = O(nlogn)，最坏情况下建堆复杂度是O(nlogn)。\\
(2)最坏情况下，树上所有的元素都要比较至叶结点。所以第i层(0<=i<=logn，根\\
为第0层)的结点最多需要比较logn-i次。而第i层有2i个结点，因此总的比较次数\\
为\\
所以最坏情况下建堆复杂度是O(n)。\\
(3)筛选法要快于本题的算法。这是因为筛选法建堆越往下层的结点比较次数越\\
少，第i层的结点最多比较logn-i次；而本题的算法越往下层的结点比较次数越多，\\
第i层的结点最多要比较i次。但是对于二叉树而言，“下层”的结点要远远多于\\
“上层”的结点，所以本题算法的比较次数要多于筛选法。\\
\\
\textbackslash{}clearpage\\


\end{solution}
\end{question}
